% ============================================================
%  Глава 12 — Реализация: MATLAB → YAML → Python
% ============================================================
\chapter{Реализация: путь чисел}
\label{ch:implementation}

\section{Вопрос}

Все предыдущие главы говорили: «матрица $\mat{A}$ имеет вот такой
вид», «$\mat{Q}$ выбирают так-то». Но как конкретное число $-6.667$ из
\eqref{eq:AB_numeric} попадает в Python-код, который шлёт команды
мотору? Где это происходит, и что меняется на каждом шаге?

\begin{ideabox}
В проекте --- классическая трёхэтапная схема:
\begin{enumerate}
  \item \textbf{Offline в MATLAB}: синтез матриц (модель, веса,
        проверка устойчивости).
  \item \textbf{Экспорт в YAML}: матрицы пишутся в \filepath{config.yaml}
        как обычные списки чисел.
  \item \textbf{Online в Python}: \texttt{mps\_node} читает YAML при старте,
        строит \texttt{StateSpaceModel} и \texttt{MPCController}, гонит
        тики.
\end{enumerate}
Плюс \textbf{live-канал} от UI: пользователь меняет матрицы в браузере
$\to$ MQTT \texttt{mps/matrices/set} $\to$ \texttt{mpc.rebuild()} на лету.
\end{ideabox}

\section{Этап 1: MATLAB (\filepath{matlab/main.m})}

В \filepath{matlab/} лежат оффлайн-скрипты для синтеза матриц.
Логика стандартная:

\begin{enumerate}
  \item Загрузить параметры объекта ($v_0, \tau_v, \tau_\omega$) из конфига.
  \item Построить $\mat{A}, \mat{B}$ по формулам \eqref{eq:A_continuous},
        \eqref{eq:B_continuous}.
  \item Проверить управляемость: \texttt{rank(ctrb(A, B)) == n}.
  \item Выбрать $\mat{Q}, \mat{R}$.
  \item Решить DARE (MATLAB-функция \texttt{dlqr}/\texttt{dare}) ---
        получить $\mat{P}, \mat{K}$.
  \item Визуализировать step response: убедиться, что замкнутая
        система устойчива и качественно ведёт себя.
  \item Экспортировать всё в YAML.
\end{enumerate}

\begin{codebox}[title={\filepath{matlab/main.m} (упрощённый эскиз)}]
\begin{lstlisting}[style=matlabstyle]
% Параметры объекта
v0    = 0.20;
tau_v = 0.15;
tau_w = 0.10;
Ts    = 0.02;     % период МПС-цикла, 50 Гц

% Непрерывные матрицы
A = [0 0 0    1         0;
     0 0 v0   0         0;
     0 0 0    0         1;
     0 0 0   -1/tau_v   0;
     0 0 0    0        -1/tau_w];
B = [0 0; 0 0; 0 0; 1/tau_v 0; 0 1/tau_w];

% Управляемость
assert(rank(ctrb(A, B)) == 5);

% Дискретизация
sysd  = c2d(ss(A, B, eye(5), zeros(5,2)), Ts, 'zoh');
Ad = sysd.A;  Bd = sysd.B;

% Веса
Q = diag([10 10 5 1 1]);
R = diag([1 1]);

% LQR / DARE
[K, P, ~] = dlqr(Ad, Bd, Q, R);

% Проверка устойчивости замкнутой
assert(all(abs(eig(Ad - Bd*K)) < 1));

% Экспорт
export_to_yaml(A, B, Q, R, K, P, 'config.yaml');
\end{lstlisting}
\end{codebox}

\begin{notebox}
MATLAB используется потому что: (1) исторически он стандарт ТАУ-курсов,
(2) \texttt{c2d}, \texttt{dlqr}, \texttt{step}, графики --- встроенные
и предсказуемые, (3) учебник Козлова пишет в MATLAB-нотации.

Альтернатива --- Python+\texttt{scipy.signal}+\texttt{matplotlib}. Сейчас
проект использует MATLAB для синтеза и Python --- для онлайн-управления.
В \filepath{tools/} есть Python-эквивалент для тех, у кого нет лицензии
MATLAB.
\end{notebox}

\section{Этап 2: YAML --- источник правды}

\filepath{config.yaml} --- общая «истина» для Python и UI. В нём
\emph{две независимые} секции для двух моделей:

\begin{codebox}[title={\filepath{config.yaml}, секция \texttt{control} (легаси)}]
\begin{verbatim}
control:
  plant:
    Ts:     0.05
    v0:     0.20
    tau_v:  0.15
    tau_w:  0.10
  matrices:
    A: [[0, 0, 0,    1,         0],
        [0, 0, 0.2,  0,         0],
        [0, 0, 0,    0,         1],
        [0, 0, 0,   -6.667,     0],
        [0, 0, 0,    0,        -10]]
    B: [[0,      0],
        [0,      0],
        [0,      0],
        [6.667,  0],
        [0,     10]]
    K_mpc: [...]    # предвычисленное усиление
    Pf:    [...]    # терминальный штраф
  weights:
    Q_diag: [10.0, 10.0, 5.0, 1.0, 1.0]
    R_diag: [1.0, 1.0]
  mpc:
    horizon_N: 10
    solver:    "clip"
  limits:
    u_min: [-0.30, -2.0]
    u_max: [+0.30, +2.0]
\end{verbatim}
\end{codebox}

\begin{codebox}[title={\filepath{config.yaml}, секция \texttt{mps} (учебная)}]
\begin{verbatim}
mps:
  plant:
    Ts: 0.02       # 50 Гц
  matrices:
    A: [...]      # непрерывная (контракт api.md)
    B: [...]      # непрерывная
    C: [...]      # для UI визуализации (по умолч. I)
    D: [...]      # для UI визуализации (по умолч. 0)
  weights:
    Q_diag: [...]
    R_diag: [...]
  horizon_N: 20
  limits:
    u_min: [...]
    u_max: [...]
  scenario:
    distance_max:           5.0
    v_target_max:           0.30
    omega_max_in_forward:   0.5
    omega_max_in_turn:      1.0
    turn_tolerance_rad:     0.05
    turn_timeout_s:         10.0
    default_distance:       2.0
    default_v_target:       0.15
\end{verbatim}
\end{codebox}

\begin{warnbox}
\textbf{Две секции не пересекаются.} \texttt{control.*} --- для старого
PID-/легаси-MPC контура; \texttt{mps.*} --- для учебного МПС-модуля.
\emph{Не путать}: легаси-секция содержит \emph{дискретные} матрицы и
готовое усиление \texttt{K\_mpc}, а МПС-секция --- \emph{непрерывные}
матрицы, и ZOH-дискретизация делается в Python при загрузке.

Если случайно подать \texttt{mps}-матрицы в легаси-контур (или наоборот),
размерности совпадут (по $5\times 5$), но смысл компонент состояния --- разный,
и MPC «закрутит» робота не туда. Так и было в коммите \texttt{8521742}
до его починки.
\end{warnbox}

\section{Этап 3: Python (\texttt{StateSpaceModel}, \texttt{MPCController})}

При старте \texttt{mps\_node} читает \texttt{config.yaml} и собирает
plant + MPC:

\begin{codebox}[title={\filepath{pi\_nodes/nodes/mps\_node.py:180}, метод \texttt{\_build\_from\_config}}]
\begin{lstlisting}[language=Python]
def _build_from_config(self):
    A = self._cfg('mps.matrices.A', None)
    B = self._cfg('mps.matrices.B', None)
    C = self._cfg('mps.matrices.C', None)
    D = self._cfg('mps.matrices.D', None)
    Q_diag = self._cfg('mps.weights.Q_diag', None)
    R_diag = self._cfg('mps.weights.R_diag', None)
    N = self._cfg('mps.horizon_N', None)
    u_min = self._cfg('mps.limits.u_min', None)
    u_max = self._cfg('mps.limits.u_max', None)
    if any(v is None for v in (A, B, Q_diag, R_diag, N, u_min, u_max)):
        raise RuntimeError('config.yaml: секция mps: неполна')

    A = np.asarray(A, dtype=float)
    B = np.asarray(B, dtype=float)
    Ad, Bd = zoh_discretize(A, B, self._mps_ts)
    plant = StateSpaceModel(Ad=Ad, Bd=Bd, Cd=C, Dd=D, Ts=self._mps_ts)
    mpc = MPCController(
        Ad=Ad, Bd=Bd,
        Q=np.diag(np.asarray(Q_diag, dtype=float)),
        R=np.diag(np.asarray(R_diag, dtype=float)),
        N=int(N),
        u_min=np.asarray(u_min, dtype=float),
        u_max=np.asarray(u_max, dtype=float),
        solver='clip',
    )
    return plant, mpc
\end{lstlisting}
\end{codebox}

Что важно отметить:

\begin{itemize}
  \item Матрицы из YAML --- \emph{непрерывные}; ZOH-дискретизация
        происходит \emph{здесь}, в одной точке кода.
  \item \texttt{Ad, Bd} затем используются \emph{и} в plant
        (\texttt{StateSpaceModel}), \emph{и} в MPC --- единый источник
        правды. Симулятор на ноуте делает то же самое (\texttt{compute\_node/
        mps\_runner.py: \_build\_controller}) --- три места, одна логика.
  \item \texttt{solver='clip'} зашит --- учебный, быстрый, не требует
        scipy на Pi.
\end{itemize}

\section{Live Apply: четвёртый канал}

Сверх старта-из-конфига есть \textbf{live-канал}: пользователь
редактирует $\mat{A}, \mat{B}, \mat{Q}, \mat{R}$ в UI и жмёт Apply.

\subsection{Поток данных}

\begin{center}
\begin{tikzpicture}[
  node distance=8mm,
  pipestep/.style={
    draw, rounded corners=2pt, fill=blue!8,
    align=center, font=\footnotesize, minimum width=4.5cm, minimum height=8mm,
  },
  flow/.style={-{Stealth[length=2mm]}, thick},
]
  \node[pipestep] (s1) {UI: \texttt{MatrixEditor} (React)};
  \node[pipestep, below=of s1] (s2) {POST \texttt{/api/v1/mps/matrices/apply}};
  \node[pipestep, below=of s2] (s3) {\filepath{routers/mps.py}\\publish MQTT};
  \node[pipestep, below=of s3] (s4) {Pi: \texttt{mps\_node.\_on\_matrices\_set}};
  \node[pipestep, below=of s4] (s5) {\texttt{plant.reload()} + \texttt{mpc.rebuild()}};
  \node[pipestep, below=of s5, fill=green!8] (s6) {publish \texttt{mps/matrices/applied}};
  \node[pipestep, below=of s6] (s7) {UI: \texttt{state.mps.applied = matrices}};
  \foreach \a/\b in {s1/s2, s2/s3, s3/s4, s4/s5, s5/s6, s6/s7}
    \draw[flow] (\a) -- (\b);
\end{tikzpicture}
\end{center}

Ключевые моменты:

\begin{itemize}
  \item Запрос отбрасывается, если прогон идёт (\texttt{is\_running ==
        True}). Apply --- \emph{между} прогонами.
  \item На Pi A, B приходят в непрерывной форме --- ZOH-дискретизируем
        тем же \texttt{zoh\_discretize}, что и на старте.
  \item \texttt{rebuild()} --- атомарный: snapshot всех полей, при
        ошибке валидации (например, $\mat{R}$ не положительная) ---
        rollback к снимку. См. главу~\ref{ch:mpc}, §rebuild.
  \item Ack \texttt{mps/matrices/applied} приходит обратно в UI ---
        пользователь видит «зелёную галочку».
\end{itemize}

\subsection{Что валидирует Pi}

\begin{enumerate}
  \item Формы матриц: $\mat{A}$ квадратная, $\mat{B}$ совместима с
        $\mat{A}$, $\mat{Q}$ диагональная нужного размера, и т.д.
  \item Положительность: $r_i > 0$. Иначе DARE не сойдётся.
  \item Размер $N \ge 1$, $u_\text{min} < u_\text{max}$.
  \item Замкнутая система устойчива (опционально, в \texttt{mpc.is\_stable()}).
\end{enumerate}

При нарушении публикуется \texttt{mps/error}, тип
\texttt{error\_type="precondition"}. UI показывает алерт.

\section{Согласование sim и робота}

\begin{ideabox}
\textbf{Кошмарный сценарий:} симулятор и робот строят MPC \emph{по-разному}.
В sim матрицы дискретизируются с $T_s = 0.02$, на роботе --- $0.05$.
Или симулятор берёт $Q$ из YAML, а робот --- свои дефолты. Тогда
\emph{прогон в sim} перестаёт быть предсказанием реального прогона ---
учебный модуль теряет смысл.

\textbf{Решение:} \emph{три} места, где строится контроллер
(на старте Pi, при Apply на Pi, в симуляторе), должны быть \textbf{строго
идентичны} по логике. Это инвариант, который тесты охраняют
(\filepath{tests/test\_mps\_runner.py}, \filepath{tests/test\_mps\_node.py}).
\end{ideabox}

\section{Тесты}

\begin{center}
\begin{tabular}{lll}
\toprule
\textbf{Файл} & \textbf{Что проверяет} & \textbf{Где живёт}\\
\midrule
\texttt{test\_mps\_runner.py}     & Sim: propagation, reach, метрики & \filepath{tests/}\\
\texttt{test\_mps\_router.py}     & REST/WS контракты, history       & \filepath{tests/}\\
\texttt{test\_mps\_node.py}       & MQTT handlers, FSM, abort        & \filepath{tests/}\\
\texttt{test\_state\_space\_extended.py} & Cd, Dd, reload(), output() & \filepath{tests/}\\
\texttt{test\_mpc\_controller.py} & rebuild(), горизонт, веса, rollback & \filepath{tests/}\\
\texttt{targetAngle.test.ts}      & Математика пикера (jsdom)        & \filepath{compute\_node/frontend/}\\
\texttt{MpsPage.test.tsx}         & UI flow: пикер, Run, история     & \filepath{compute\_node/frontend/}\\
\bottomrule
\end{tabular}
\end{center}

Coverage таргет --- 85\% по проекту. CI на каждый PR в feat/mps
прогоняет все эти суиты, плюс \emph{smoke-тест на реальном Pi} (A10)
--- руками, по чек-листу.

\section{Где какие риски}

\begin{notebox}
Учебная природа модуля = вольный экспериментатор у руля. Места, где
\emph{реально} может «не пойти», и куда смотреть в первую очередь:

\begin{itemize}
  \item \textbf{$\mat{R}$ с нулём.} DARE не сойдётся, \texttt{rebuild()}
        упадёт. Pi ответит \texttt{precondition}.
  \item \textbf{$\mat{Q}$ с большими разбросами весов.} Может дать
        численно плохо обусловленную $\mat{H}$, $\mat{K}$ будет
        огромным $\to$ робот будет дёргаться. Pre-validate этого
        не ловит, ловит \texttt{is\_stable()}.
  \item \textbf{Несовпадение секций конфига и трактовки.} См. warning
        выше про коммит \texttt{8521742}.
  \item \textbf{Watchdog по одометрии.} Если выдернуть USB у Pi от
        мотор-контроллера во время прогона --- \texttt{watchdog}
        отработает за $3 \cdot T_s = 60$~мс.
\end{itemize}
\end{notebox}

\section{Что мы получили}

\begin{itemize}
  \item Числа модели проходят через три точки: MATLAB (синтез) ---
        YAML (источник правды) --- Python (онлайн).
  \item Кроме старта-из-конфига, есть live-канал: UI \texttt{->} MQTT
        \texttt{->} \texttt{mpc.rebuild()}. Между прогонами, атомарно,
        с rollback.
  \item Три места построения контроллера (Pi-start, Pi-Apply, simulator)
        обязаны быть идентичны --- это инвариант.
  \item Тестов --- семь файлов, покрывающих все три уровня (юнит на
        backend, юнит на frontend, smoke на Pi).
\end{itemize}

\medskip
\textbf{Поздравляем.} От «нажал кнопку» до ШИМ-сигнала в моторе ---
вся математика, вся архитектура, все подводные камни. Если ты
читал внимательно, ты теперь можешь:

\begin{enumerate}
  \item Объяснить другому, что значит «вектор состояния».
  \item Вывести $\mat{A}, \mat{B}$ из физической постановки.
  \item Понимать, что делает каждая строка \texttt{mps\_node.\_tick()}.
  \item Предсказать, как изменится поведение робота при изменении
        $\mat{Q}, \mat{R}, N$.
  \item Найти и починить баг, аналогичный коммиту \texttt{8521742}.
\end{enumerate}

\textbf{Удачи с курсовой.}
